\documentclass[a4paper, 12pt]{mksheetphhd}

\title{Programmierübungen}
\author{Marvin Ködding}
\usepackage{minted}
\usepackage{multicol}

\everymath{\displaystyle}

\begin{document}
	\printtitle
	
	\vspace*{.5cm}
	
	\begin{aufgabe}
		Ein Primzahlzwilling ist ein Paar von Primzahlen, deren Differenz zwei ist. Der erste Primzahlzwilling
		ist $(3, 5)$. Welches ist der siebte Primzahlzwilling?
	\end{aufgabe}
	
	\begin{aufgabe}
		Ein Primzahlzwilling ist ein Paar von Primzahlen, deren Differenz zwei ist. Wie viele Primzahlzwillinge
		liegen zwischen $2000$ und $4000$?
	\end{aufgabe}
	
	\begin{aufgabe}
		Wie viele verschiedene Zeichenketten der Länge $4$, die mindestens eine der folgenden Bedingungen erfüllen, kann man aus den $8$ Buchstaben \texttt{A} bis \texttt{H} konstruieren?
		\begin{enumerate}
			\item Der erste Buchstabe ist \texttt{A}.
			\item Der zweite Buchstabe ist \texttt{B}.
			\item Der dritte Buchstabe ist \texttt{C}.
		\end{enumerate}
		Erlaubt sind also z.B. die Zeichenketten \texttt{ADBH} oder \texttt{DBDH}, aber auch \texttt{AHCC} oder \texttt{ABCA}.
	\end{aufgabe}
	
	\begin{aufgabe}
		Eine Dreieckszahl ist eine positive Zahl, die sich als Summe der Form $\sum_{k=1}^n k$ darstellen lässt. Die ersten fünf Dreieckszahlen sind also $1$, $3$, $6$, $10$ und $15$. Wie viele der ersten $500$ Dreieckszahlen sind Quadratzahlen? Und wie viele der ersten hundert positiven Quadratzahlen sind Dreieckszahlen?
	\end{aufgabe}
	
	\begin{aufgabe}
		Gebt die kleinste Fibonacci-Zahl an, die größer als eins und ein Quadrat einer natürlichen Zahl ist.
	\end{aufgabe}
	
	\printlicense
	
	\printsocials
	
\end{document}